%MIT OpenCourseWare: https://ocw.mit.edu
%18.100A / 18.1001 Real Analysis, Fall 2020
%License: Creative Commons BY-NC-SA 
%For information about citing these materials or our Terms of Use, visit: https://ocw.mit.edu/terms.

\documentclass{article}
\usepackage[utf8]{inputenc}
\usepackage{pkige}
\usepackage{wrapfig}

\title{18.100A: Complete Lecture Notes}

%\title{18.100A: Typed Recitation Notes}
%%%\author{Recitation 01}
%%\author{Recitation 02}
%%\author{Recitation 03}
%%\author{Recitation 04}
%%\author{Recitation 05}
%%\author{Recitation 06}
%%\author{Midterm Exam Review}
%%\author{Final Exam Review}

%\author{Lecture 1: \\ Sets, Set Operations, and Mathematical Induction}
%\author{Lecture 2: \\ Cantor's Theory of Cardinality (Size)}
%\author{Lecture 3: \\ Cantor's Remarkable Theorem \\ and the Rationals' Lack of the Least Upper Bound Property}
%\author{Lecture 4: \\ 	The Characterization of the Real Numbers}
%\author{Lecture 5: \\ 	The Archimedian Property, Density of the Rationals, and Absolute Value}
%\author{Lecture 6: \\ 	The Uncountabality of the Real Numbers}
%\author{Lecture 7: \\ 	Convergent Sequences of Real Numbers}
%\author{Lecture 8: \\ 	The Squeeze Theorem and Operations Involving Convergent Sequences}
%\author{Lecture 9: \\ Limsup, Liminf, and the Bolzano-Weierstrass Theorem}
%\author{Lecture 10: \\ The Completeness of the Real Numbers and Basic Properties of Infinite Series}
%\author{Lecture 11: \\ Absolute Convergence and the Comparison Test for Series}
%\author{Lecture 12: \\ The Ratio, Root, and Alternating Series Tests}
%\author{Lecture 13: \\ Limits of Functions}
%\author{Lecture 14: \\ Limits of Functions in Terms of Sequences and Continuity}
%\author{Lecture 15: \\ The Continuity of Sine and Cosine \\ and the Many Discontinuities of Dirichlet's Function}
%\author{Lecture 16: \\ The Min/Max Theorem and Bolzano's Intermediate Value Theorem}
%\author{Lecture 17: \\ Uniform Continuity and the Definition of the Derivative}
%\author{Lecture 18: \\ 	Weierstrass's Example of a Continuous and Nowhere Differentiable Function}
%\author{Lecture 19: \\ Differentiation Rules, Rolle's Theorem, and the Mean Value Theorem}
%\author{Lecture 20: \\ Taylor's Theorem and the Definition of Riemann Sums}
%\author{Lecture 21: \\ The Riemann Integral of a Continuous Function}
%\author{Lecture 22: \\ The Fundamental Theorem of Calculus, Integration by Parts, and Change of Variable Formula}
%\author{Lecture 23: \\ 	Pointwise and Uniform Convergence of Sequences of Functions}
%\author{Lecture 24: \\ Uniform Convergence, the Weierstrass M-Test, and Interchanging Limits}
%\author{Lecture 25: \\ Power Series and the Weierstrass Approximation Theorem}

\date{}

\begin{document}
\maketitle

%---------------------------------------------------

%\input{Lectures/Lec01} \newpage 
%\input{Lectures/Lec02} \newpage 
%\input{Lectures/Lec03} \newpage
%\input{Lectures/Lec04} \newpage
%\input{Lectures/Lec05} \newpage
%\input{Lectures/Lec06} \newpage
%\input{Lectures/Lec07} \newpage
%\input{Lectures/Lec08} \newpage
%\input{Lectures/Lec09} \newpage
%\input{Lectures/Lec10} \newpage
%\input{Lectures/Lec11} \newpage
%\input{Lectures/Lec12} \newpage
%\input{Lectures/Lec13} \newpage
%\input{Lectures/Lec14} \newpage
%\input{Lectures/Lec15} \newpage
%\input{Lectures/Lec16} \newpage
%\input{Lectures/Lec17} \newpage
%\input{Lectures/Lec18} \newpage
%\input{Lectures/Lec19} \newpage
%\input{Lectures/Lec20} \newpage
%\input{Lectures/Lec21} \newpage
%\input{Lectures/Lec22} \newpage
%\input{Lectures/Lec23} \newpage
%\input{Lectures/Lec24} \newpage
%\input{Lectures/Lec25} \newpage

%---------------------------------------------------

%%\section*{Recitation Notes}

%\input{Recitations/Rec01}\newpage
%\input{Recitations/Rec02}\newpage
%\input{Recitations/Rec03}\newpage
%\input{Recitations/Rec04}\newpage
%\input{Recitations/Rec05}\newpage
%\input{Recitations/Rec06}\newpage
%\input{Recitations/Rec07_Midterm}\newpage 
%\input{Recitations/Rec08_Final} \newpage


%---------------------------------------------------

%%\section*{Lecture Notes} WITH TOC

\tableofcontents
\newpage
   
\section[ Sets, Set Operations, and Mathematical Induction]{Lecture 1 \\ {\normalsize Sets, Set Operations, and Mathematical Induction}}\input{Lectures/Lec01} \newpage 
\section[ Cantor's Theory of Cardinality (Size)]{Lecture 2 \\ {\normalsize Cantor's Theory of Cardinality (Size)}} \input{Lectures/Lec02} \newpage 
\section[Cantor's Remarkable Theorem and the Rationals' Lack of the Least Upper Bound Property]{Lecture 3 \\ {\normalsize Cantor's Remarkable Theorem and the Rationals' Lack of the Least Upper Bound Property}} \input{Lectures/Lec03} \newpage 
\section[The Characterization of the Real Numbers]{Lecture 4 \\ {\normalsize The Characterization of the Real Numbers}}\input{Lectures/Lec04} \newpage 
\section[The Archimedian Property, Density of the Rationals, and Absolute Value]{Lecture 5 \\ {\normalsize The Archimedian Property, Density of the Rationals, and Absolute Value}} \input{Lectures/Lec05} \newpage 
\section[The Uncountabality of the Real Numbers]{Lecture 6 \\ {\normalsize The Uncountabality of the Real Numbers}}\input{Lectures/Lec06} \newpage 
\section[Convergent Sequences of Real Numbers]{Lecture 7 \\ {\normalsize 	Convergent Sequences of Real Numbers}}\input{Lectures/Lec07} \newpage 
\section[The Squeeze Theorem and Operations Involving Convergent Sequences]{Lecture 8 \\ {\normalsize 	The Squeeze Theorem and Operations Involving Convergent Sequences}}\input{Lectures/Lec08} \newpage 
\section[Limsup, Liminf, and the Bolzano-Weierstrass Theorem]{Lecture 9 \\ {\normalsize Limsup, Liminf, and the Bolzano-Weierstrass Theorem}}\input{Lectures/Lec09} \newpage 
\section[The Completeness of the Real Numbers and Basic Properties of Infinite Series]{Lecture 10 \\ {\normalsize The Completeness of the Real Numbers and Basic Properties of Infinite Series}} \input{Lectures/Lec10} \newpage 
\section[Absolute Convergence and the Comparison Test for Series]{Lecture 11 \\ {\normalsize Absolute Convergence and the Comparison Test for Series}}\input{Lectures/Lec11} \newpage 
\section[The Ratio, Root, and Alternating Series Tests]{Lecture 12 \\ {\normalsize The Ratio, Root, and Alternating Series Tests}}\input{Lectures/Lec12} \newpage 
\section[Limits of Functions]{Lecture 13 \\ {\normalsize Limits of Functions}} \input{Lectures/Lec13} \newpage 
\section[Limits of Functions in Terms of Sequences and Continuity]{Lecture 14 \\ {\normalsize Limits of Functions in Terms of Sequences and Continuity}}\input{Lectures/Lec14} \newpage 
\section[The Continuity of Sine and Cosine and the Many Discontinuities of Dirichlet's Function]{Lecture 15 \\ {\normalsize The Continuity of Sine and Cosine and the Many Discontinuities of Dirichlet's Function}}\input{Lectures/Lec15} \newpage 
\section[The Min/Max Theorem and Bolzano's Intermediate Value Theorem]{Lecture 16 \\ {\normalsize The Min/Max Theorem and Bolzano's Intermediate Value Theorem}}\input{Lectures/Lec16} \newpage 
\section[Uniform Continuity and the Definition of the Derivative]{Lecture 17 \\ {\normalsize Uniform Continuity and the Definition of the Derivative}}\input{Lectures/Lec17} \newpage 
\section[Weierstrass's Example of a Continuous and Nowhere Differentiable Function]{Lecture 18 \\ {\normalsize 	Weierstrass's Example of a Continuous and Nowhere Differentiable Function}}\input{Lectures/Lec18} \newpage 
\section[Differentiation Rules, Rolle's Theorem, and the Mean Value Theorem]{Lecture 19 \\ {\normalsize Differentiation Rules, Rolle's Theorem, and the Mean Value Theorem}}\input{Lectures/Lec19} \newpage 
\section[Taylor's Theorem and the Definition of Riemann Sums]{Lecture 20 \\ {\normalsize Taylor's Theorem and the Definition of Riemann Sums}}\input{Lectures/Lec20} \newpage 
\section[The Riemann Integral of a Continuous Function]{Lecture 21 \\ {\normalsize The Riemann Integral of a Continuous Function}}\input{Lectures/Lec21} \newpage 
\section[The Fundamental Theorem of Calculus, Integration by Parts, and Change of Variable Formula]{Lecture 22 \\ {\normalsize The Fundamental Theorem of Calculus, Integration by Parts, and Change of Variable Formula}}\input{Lectures/Lec22} \newpage 
\section[Pointwise and Uniform Convergence of Sequences of Functions]{Lecture 23 \\ {\normalsize 	Pointwise and Uniform Convergence of Sequences of Functions}}\input{Lectures/Lec23} \newpage 
\section[Uniform Convergence, the Weierstrass M-Test, and Interchanging Limits]{Lecture 24 \\ {\normalsize Uniform Convergence, the Weierstrass M-Test, and Interchanging Limits}}\input{Lectures/Lec24} \newpage 
\section[Power Series and the Weierstrass Approximation Theorem]{Lecture 25 \\ {\normalsize Power Series and the Weierstrass Approximation Theorem}}\input{Lectures/Lec25}

\end{document}